\begin{figure}[H]
\centering
\begin{tikzpicture}[scale=1.2]
% Grid World Environment
\draw[step=1cm,gray,very thin] (0,0) grid (12,8);
\draw[thick] (0,0) rectangle (12,8);

% Generation Counter
\node[rectangle,draw=black,fill=yellow!20,thick,font=\small\bfseries] at (10.5,8.5) {Gen: 50/1000};

% ===== FAMILY 1 (Top Left) =====
% Mother 1
% \node[circle,draw=red!70,fill=red!30,thick,minimum size=0.5cm,drop shadow] (m1) at (3.5,6.5) {M1};
% \node[font=\tiny] at (3.5,5.9) {Mother 1};
\node[circle,draw=red!70,fill=red!30,thick,
      minimum size=0.35cm,inner sep=0pt,drop shadow] (m1) at (3.5,6.5) {M1};
\node[font=\tiny] at (3.5,5.9) {Mother 1};


% % Child 1
\node[circle,draw=blue!70,fill=blue!30,thick,minimum size=0.5cm,inner sep=0pt,drop shadow] (c1) at (1.5,6.5) {C1};
\node[font=\tiny] at (1.5,5.9) {Child 1};


% Bond 1
\draw[line width=1.5pt,purple!60,decorate,decoration={snake,amplitude=1.5pt,segment length=6pt}] (m1) -- (c1);

% ===== FAMILY 2 (Top Right) =====

% Mother 2
\node[circle,draw=red!70,fill=red!30,thick,minimum size=0.7cm,inner sep=0pt,drop shadow] (m2) at (8,6.5) {M2};
\node[font=\tiny] at (8,5.9) {Mother 2};

% Child 2
\node[circle,draw=blue!70,fill=blue!30,thick,minimum size=0.6cm,inner sep=0pt,drop shadow] (c2) at (10.5,6.5) {C2};
\node[font=\tiny] at (10.5,5.9) {Child 2};



% Bond 2
\draw[line width=1.5pt,purple!60,decorate,decoration={snake,amplitude=1.5pt,segment length=6pt}] (m2) -- (c2);

% ===== FAMILY 3 (Bottom Center) =====
% Mother 3
\node[circle,draw=red!70,fill=red!30,thick,minimum size=0.7cm,inner sep=0pt,drop shadow] (m3) at (6,3.5) {M3};
\node[font=\tiny] at (6,2.9) {Mother 3};

% Child 3
\node[circle,draw=blue!70,fill=blue!30,thick,minimum size=0.6cm,inner sep=0pt,drop shadow] (c3) at (6,1.5) {C3};
\node[font=\tiny] at (6,0.9) {Child 3};



% Bond 3
\draw[line width=1.5pt,purple!60,decorate,decoration={snake,amplitude=1.5pt,segment length=6pt}] (m3) -- (c3);

% ===== SHARED RESOURCES =====
% Threat (shared danger)
\node[regular polygon,regular polygon sides=3,draw=black!70,fill=gray!40,thick,minimum size=0.7cm] (threat) at (6,5) {T};
\node[font=\tiny] at (6,4.4) {Threat};

% Food Resources (limited, shared)
\node[rectangle,draw=green!70,fill=green!30,thick,minimum size=0.5cm] (food1) at (2,3) {F};
\node[font=\tiny] at (2,2.5) {Food};

\node[rectangle,draw=green!70,fill=green!30,thick,minimum size=0.5cm] (food2) at (9.5,3) {F};
\node[font=\tiny] at (9.5,2.5) {Food};

% ===== INTERACTION ARROWS =====
% M1 sensing C2 (potential alloparental care decision)
\draw[->,thick,dashed,orange!70] (m1) to[bend left=20] (c2);
\node[orange!70,font=\small] at (5.5,7.2) {Help?};


% M2 moving toward food (competing with M3)
\draw[->,thick,blue!60] (m2) to[bend right=15] (food2);
\node[blue!60,font=\small] at (7.8,4.5) {Compete};


% M3 protecting all children from threat (cooperative defense)
\draw[->,thick,green!100] (m3) -- (threat);
\node[green!100,font=\small] at (5.4,4.2) {Protect};

% M1 foraging (caring for own child first)
\draw[->,thick,brown!60] (m1) to[bend right=15] (food1);
\node[brown!100,font=\small] at (1.6,4.5) {Forage};

% ===== SENSORY RANGES (showing overlap) =====
\draw[densely dotted,red!30,thick] (m1) circle (2.5cm);
\draw[densely dotted,red!30,thick] (m2) circle (2.5cm);
\draw[densely dotted,red!30,thick] (m3) circle (2.8cm);


% LEGEND
\node[rectangle,draw=gray!60,fill=gray!10,align=left,font=\tiny,inner sep=4pt] at (1.5,1.2) {
    \textbf{Legend:}\\
    M: Mother Agent\\
    C: Child Agent\\
    T: Shared Threat\\
    F: Shared Food\\
    Dotted: Sensory Range
};


\end{tikzpicture}
\caption{\textbf{Illustration of the multi-family evolutionary grid-world environment.} 
The environment is represented as a two-dimensional grid containing three mother--child pairs, shared food resources, and a common threat. 
Each mother agent occupies the environment together with its offspring and perceives nearby cues within a limited sensory range, shown by the dotted circles. 
Food entities provide energy for maternal self-maintenance, whereas the threat introduces environmental pressure that may trigger protective behavior. 
The arrows illustrate example behavioral tendencies that may arise in the simulation, including foraging, competition for resources, protection from threats, and potential caregiving toward non-kin offspring. 
Across generations, maternal behavioral parameters are inherited and selected according to offspring survival, allowing different caregiving strategies to emerge and evolve under shared ecological constraints.}
\label{fig:gridworld_multifamily}
\end{figure}